\documentclass{article}
\usepackage[UTF8]{ctex}
\usepackage{amsmath,amssymb,amsfonts}
\usepackage{enumitem}
\usepackage{geometry}
\geometry{a4paper, margin=1in}
% === 错题本排版优化（不参与题目标识） ===
\setlength{\parindent}{0pt}
\setlength{\parskip}{0.35em plus 0.10em minus 0.08em}
\setlength{\emergencystretch}{3em}
\setlist{itemsep=0.25em, topsep=0.40em, parsep=0pt}
\allowdisplaybreaks[2]
% === 错题本排版优化结束 ===
\begin{document}
\section*{哈尔滨工程大学2025年高等代数真题}
\section*{78．哈尔滨工程大学2025年高等代数试卷}
一、填空题。\\

\subsection{题目1}
1．已知 \(\mathbb{R}[x]\) 中多项式

\[
f(x)=x^{4}+2 x^{3}-x^{2}-4 x-2, g(x)=x^{4}+x^{3}-x^{2}-2 x-2 .
\]

则 \((f(x), g(x))=\)

\subsection{题目2}
2．已知矩阵 \(A=\left(\begin{array}{ccc}1 & 0 & 0 \\ 1 & 1 & 3 \\ 0 & 1 & -1\end{array}\right), B=\left(\begin{array}{lll}1 & 0 & 1 \\ 0 & 1 & 0 \\ 0 & 0 & 1\end{array}\right)\) ，矩阵 \(X\) 满足 \(A X A=B\) ，则 \(X=\) \(\_\_\_\_\)

\begin{itemize}

\subsection{题目3}
\end{itemize}

\subsection{题目3}
原题边界“3．”：令 \(V\) 为 \(\mathbb{R}\) 上由矩阵 \(A\) 的全体实系数多项式组成的线性空间，其中
\[
A=\left(\begin{array}{ccc}
1 & & \\
& \omega & \\
& & \omega^{2}
\end{array}\right), w=\frac{-1+\sqrt{3} \mathrm{i}}{2}, \text { 且 } \mathrm{i}^{2}=-1 .
\]
则 \(V\) 的维数是 \(\_\_\_\_\)．

\subsection{题目4}

\subsection{题目4}
原题边界“4．”：已知矩阵 \(A\) 的初等因子为 \(\lambda, \lambda^{2},(\lambda+1)^{3}, \lambda-1, \lambda-1\) ，则矩阵 \(A\) 的最小多项式是 \(\_\_\_\_\)．

\subsection{题目5}
5．已知3维欧氏空间中，内积在 \(a_{1}, a_{2}, a_{3}\) 下的度量矩阵为

\[
\left(\begin{array}{ccc}
1 & -1 & 1 \\
-1 & 2 & 0 \\
1 & 0 & 3
\end{array}\right)
\]

若 \(\beta_{1}=2 \alpha_{2}-\alpha_{3}\) 与 \(\beta_{2}=t a_{1}-3 a_{2}+a_{3}\) 正交，则 \(t=\) \(\_\_\_\_\)．

二．解答题。\\

\subsection{题目6}
6．设 \(p\) 奇素数，判别多项式 \(f(x)=x^{p}+p x+1\) 在 \(Q\) 上是否可约．

\subsection{题目7}
7．证明 \(n\) 阶行列式

\[
\left|\begin{array}{cccccc}
\cos \alpha & 1 & 0 & \cdots & 0 & 0 \\
1 & 2 \cos \alpha & 1 & \cdots & 0 & 0 \\
0 & 1 & 2 \cos \alpha & \cdots & 0 & 0 \\
\vdots & \vdots & \vdots & & \vdots & \vdots \\
0 & 0 & 0 & \cdots & 2 \cos \alpha & 1 \\
0 & 0 & 0 & \cdots & 1 & 2 \cos \alpha
\end{array}\right|=\cos n \alpha .
\]

\subsection{题目8}
8．讨论参数 \(\lambda, \mu\) 取什么值时，线性方程组

\[
\left\{\begin{array}{l}
x_{1}+x_{2}+x_{3}+x_{4}+x_{5}=1 \\
3 x_{1}+2 x_{2}+x_{3}+x_{4}-3 x_{5}=\lambda \\
x_{2}+2 x_{3}+2 x_{4}+6 x_{5}=3 \\
5 x_{1}+4 x_{2}+3 x_{3}+3 x_{4}-x_{5}=\mu
\end{array}\right.
\]

有解？并求解。

\subsection{题目9}
9．设 \(n\) 阶矩阵 \(A\) 满足 \(A^{2}=A\) 。证明：

\begin{itemize}
  \item[（1）] \(r(A)+r\left(E_{n}-A\right)=n\) ．
  \item[（2）] 存在可逆矩阵 \(P\) ，使得 \(A=P\left(\begin{array}{cc}E_{r} & O \\ O & O\end{array}\right) P^{-1}\) ，其中 \(r(A)=r\) ．
\end{itemize}

\subsection{题目10}
10．已知向量组

\[
\alpha_{1}=(2,0,1,3,-1)^{\mathrm{T}}, \alpha_{2}=(0,-2,1,5,-3)^{\mathrm{T}}, \beta_{1}=(1,1,0,-1,1)^{\mathrm{T}}, \beta_{2}=(1,-3,2,0,5)^{\mathrm{T}} .
\]

记 \(W_{1}=L\left(\alpha_{1}, \alpha_{2}\right), W_{2}=L\left(\beta_{1}, \beta_{2}\right)\) ，求 \(W_{1}+W_{2}\) 与 \(W_{1} \cap W_{2}\) 的维数与基。

\subsection{题目11}
11．设 \(A=\left(\begin{array}{ccc}4 & 5 & -2 \\ -2 & -2 & 1 \\ -1 & -1 & 1\end{array}\right)\) ，求 \(A\) 的若尔当标准形 \(J\) ，并求可逆矩阵 \(T\) ，使得 \(A=T J T^{-1}\) 。

\subsection{题目12}
12．二次型

\[
f\left(x_{1}, x_{2}, x_{3}\right)=X^{\mathrm{T}} A X
\]

在正交变换 \(X=Q Y\) 下的标准形为 \(y_{1}^{2}+y_{2}^{2}\) ，其中 \(X=\left(x_{1}, x_{2}, x_{3}\right)^{\mathrm{T}}, Y=\left(y_{1}, y_{2}, y_{3}\right)^{\mathrm{T}}\) ，且正交矩阵\(Q\) 的第三列为 \(\left(\frac{\sqrt{2}}{2}, 0, \frac{\sqrt{2}}{2}\right)^{\mathrm{T}}\) ，求实对称矩阵 \(A\) ，并证明 \(A+E_{3}\) 正定。

\subsection{题目13}
13．取 \(E_{11}, E_{12}, E_{21}, E_{22}\) 为 \(\mathbb{R}^{2 \times 2}\) 的一组基，定义线性变换

\[
\sigma(X)=\left(\begin{array}{ll}
a & b \\
c & d
\end{array}\right) X, X \in \mathbb{R}^{2 \times 2}
\]

\begin{itemize}
  \item[（1）] 求 \(\sigma\) 在基 \(E_{11}, E_{12}, E_{21}, E_{22}\) 下的矩阵。
  \item[（2）] 证明：当 \(a d-b c \neq 0\) 时，\(\sigma\) 可逆．
  \item[（3）] 当 \(\sigma\) 可逆时，求 \(\sigma^{-1}\) 在基 \(E_{11}, E_{12}, E_{21}, E_{22}\) 下的矩阵．
\end{itemize}

\subsection{题目14}
14．设 \(V\) 是 \(\mathbb{C}\) 上的 \(n\) 维线性空间，\(V\) 上线性变换 \(\tau\) 在基 \(\varepsilon_{1}, \varepsilon_{2}, \cdots, \varepsilon_{n}\) 下的矩阵

\[
A=\left(\begin{array}{llll}
\lambda & & & \\
1 & \lambda & & \\
& & \ddots & \\
& & 1 & \lambda
\end{array}\right) .
\]

证明：对任意非零的 \(\tau\)－子空间 \(W\) ，有 \(\varepsilon_{n} \in W\) 。
\end{document}
